\documentclass{article}
\usepackage{amsmath,amssymb,amsthm}
\usepackage{algorithm,algorithmic}
\usepackage{hyperref}
\newtheorem{theorem}{Theorem}
\newtheorem{lemma}{Lemma}
\newcommand{\E}{\mathbb{E}}

\begin{document}

\section*{Accelerated Gradient Method for Strongly Convex Smooth Functions (AGM-SC)}

--------------------------------------------------------------------
\section*{Mathematical Formulation}
Let $f:\mathbb{R}^{d}\rightarrow\mathbb{R}$ be \emph{$\mu$‑strongly convex} and \emph{$L$‑smooth}:
\[
\frac{\mu}{2}\|u-v\|^{2}\le f(u)-f(v)-\langle\nabla f(v),u-v\rangle\le\frac{L}{2}\|u-v\|^{2},
\qquad\forall u,v\in\mathbb{R}^{d}.
\]

Define the momentum coefficient
\[
\beta\;:=\;\frac{\sqrt{L}-\sqrt{\mu}}{\sqrt{L}+\sqrt{\mu}}\in[0,1).
\]

Given a stepsize $\alpha:=\frac{1}{L}$, the algorithm maintains two sequences
$\{x_k\}_{k\ge 0}$ and $\{y_k\}_{k\ge 0}$:
\[
\boxed{
\begin{aligned}
y_k &= x_k + \beta\,(x_k-x_{k-1}),\qquad k\ge 1,\\[4pt]
x_{k+1} &= y_k - \alpha\,\nabla f(y_k),\qquad k\ge 0,
\end{aligned}}
\]
with an initial point $x_0\in\mathbb{R}^{d}$ and a dummy previous iterate $x_{-1}=x_0$.

--------------------------------------------------------------------
\section*{Pseudocode}
\begin{algorithm}[H]
\caption{AGM‑SC}
\begin{algorithmic}[1]
\Require $x_0\in\mathbb{R}^{d}$, stepsize $\alpha=1/L$, momentum $\beta=\frac{\sqrt L-\sqrt\mu}{\sqrt L+\sqrt\mu}$, max iter $T$.
\State $x_{-1}\gets x_0$ \Comment{initialize previous iterate}
\For{$k=0,\dots,T-1$}
    \If{$k=0$}
        \State $y_k \gets x_k$ \Comment{no momentum in the first step}
    \Else
        \State $y_k \gets x_k + \beta\,(x_k-x_{k-1})$ \Comment{(1)}
    \EndIf
    \State $g_k \gets \nabla f(y_k)$
    \State $x_{k+1} \gets y_k - \alpha\, g_k$ \Comment{(2)}
\EndFor
\State \Return $x_T$
\end{algorithmic}
\end{algorithm}

--------------------------------------------------------------------
\section*{Dry Run Example}
Consider the one‑dimensional quadratic
\[
f(w)=(w-3)^{2},\qquad \nabla f(w)=2(w-3).
\]
For a quadratic $f$ we have $L=\mu=2$, hence
\[
\alpha=\frac{1}{L}=0.5,\qquad 
\beta=\frac{\sqrt{L}-\sqrt{\mu}}{\sqrt{L}+\sqrt{\mu}}=0.
\]
Thus the algorithm reduces to plain gradient descent.

\begin{center}
\begin{tabular}{c|c|c|c}
iteration $k$ & $y_k$ & $x_{k+1}=y_k-\alpha\nabla f(y_k)$ & $f(x_{k+1})$\\ \hline
0 & $y_0=x_0=0$ &
$0-0.5\cdot 2(0-3)=3$ & $(3-3)^2=0$\\ \hline
1 & $y_1=x_1+0\cdot(x_1-x_0)=3$ &
$3-0.5\cdot 2(3-3)=3$ & $0$\\ \hline
2 & $y_2=3$ &
$3-0.5\cdot 2(3-3)=3$ & $0$
\end{tabular}
\end{center}
The method reaches the optimum $w^{\ast}=3$ after a single iteration, illustrating the linear convergence guarantee for strongly convex quadratics.

--------------------------------------------------------------------
\section*{Assumptions}
\begin{enumerate}
\item[(A1)] \textbf{Strong convexity}: $f$ is $\mu$‑strongly convex with $\mu>0$.
\item[(A2)] \textbf{Smoothness}: $f$ has $L$‑Lipschitz gradient with $L\ge\mu$.
\item[(A3)] \textbf{Exact gradients}: at each iteration the true gradient $\nabla f(y_k)$ is available.
\item[(A4)] \textbf{Parameter choice}: stepsize $\alpha=1/L$ and momentum $\beta$ as defined above.
\end{enumerate}
All constants $L,\mu$ are known (or estimated) beforehand.

--------------------------------------------------------------------
\section*{Main Convergence Theorem}
\begin{theorem}[Linear convergence of AGM‑SC]\label{thm:linear}
Assume (A1)–(A4). Let $x^{\ast}=\arg\min f$ and define $q:=1-\sqrt{\mu/L}\in(0,1)$. Then for all $k\ge 0$,
\[
\boxed{f(x_k)-f(x^{\ast})\;\le\; q^{\,k}\,\bigl(f(x_0)-f(x^{\ast})\bigr)}.
\tag{3}
\]
Equivalently,
\[
\|x_k-x^{\ast}\|^{2}\le \frac{2}{\mu}\,q^{\,k}\,\bigl(f(x_0)-f(x^{\ast})\bigr).
\tag{4}
\]
\end{theorem}
Thus AGM‑SC attains the optimal accelerated linear rate $1-\sqrt{\mu/L}$.

--------------------------------------------------------------------
\section*{Theoretical Properties}
\begin{itemize}
\item \textbf{Stability}: The Lyapunov function constructed in the proof is non‑increasing for any $\beta$ in $[0,1)$; the specific choice (A4) yields the smallest contraction factor $q$.
\item \textbf{Hyper‑parameter sensitivity}: If $\beta$ deviates from the optimal value, the rate becomes $1-\Theta(\beta\sqrt{\mu/L})$; the algorithm remains convergent as long as $\beta<1$.
\item \textbf{Comparison to baselines}:
    \begin{itemize}
        \item Gradient Descent (GD) with $\alpha=1/L$ enjoys rate $1-\mu/L$.
        \item AGM‑SC improves the factor from $1-\mu/L$ to $1-\sqrt{\mu/L}$, which is provably optimal for first‑order methods on $\mu$‑strongly convex $L$‑smooth functions.
    \end{itemize}
\end{itemize}

--------------------------------------------------------------------
\section*{Discussion}
The theorem shows that the accelerated scheme reduces the optimality gap exponentially fast, with a contraction factor that depends on the square‑root of the condition number $\kappa:=L/\mu$. This matches the lower bound of Nesterov (1983) and explains the empirical superiority of momentum‑based methods on ill‑conditioned problems.

--------------------------------------------------------------------
\section*{Preliminaries (Definitions and Lemmas)}
\begin{lemma}[Smoothness inequality]\label{lem:smooth}
For an $L$‑smooth function,
\[
f(u)\le f(v)+\langle\nabla f(v),u-v\rangle+\frac{L}{2}\|u-v\|^{2},\qquad\forall u,v.
\tag{5}
\]
\end{lemma}
\begin{lemma}[Strong convexity inequality]\label{lem:strong}
For a $\mu$‑strongly convex function,
\[
f(u)\ge f(v)+\langle\nabla f(v),u-v\rangle+\frac{\mu}{2}\|u-v\|^{2},\qquad\forall u,v.
\tag{6}
\]
\end{lemma}
\begin{lemma}[Three‑point identity]\label{lem:three}
For any $a,b,c\in\mathbb{R}^{d}$,
\[
2\langle a-b,c-b\rangle=\|a-c\|^{2}-\|a-b\|^{2}-\|c-b\|^{2}.
\tag{7}
\]
\end{lemma}

--------------------------------------------------------------------
\section*{Main Proof (Step‑by‑Step Derivation)}
We prove Theorem~\ref{thm:linear} by constructing a Lyapunov function and showing it contracts by factor $q$ each iteration.

\noindent\textbf{Step 1 (Define potential).}  
Let $x^{\ast}$ be the unique minimizer of $f$. Set
\[
\Phi_k\;:=\; (1+\gamma)\bigl(f(x_k)-f(x^{\ast})\bigr)+\frac{1}{2\alpha}\|x_k-x^{\ast}\|^{2},
\tag{8}
\]
where $\gamma:=\frac{1-\sqrt{\mu/L}}{1+\sqrt{\mu/L}}=\beta$ (the same as the momentum).  
The choice of $\gamma$ will make the cross terms cancel.

\noindent\textbf{Step 2 (Relation between $x_{k+1}$ and $y_k$).}  
From the update $x_{k+1}=y_k-\alpha\nabla f(y_k)$ we have
\[
\|x_{k+1}-x^{\ast}\|^{2}
= \|y_k-x^{\ast}\|^{2}-2\alpha\langle\nabla f(y_k),y_k-x^{\ast}\rangle+\alpha^{2}\|\nabla f(y_k)\|^{2}.
\tag{9}
\]

\noindent\textbf{Step 3 (Apply smoothness to bound $\|\nabla f(y_k)\|^{2}$).}  
By Lemma~\ref{lem:smooth} with $u=y_k-\alpha\nabla f(y_k)=x_{k+1}$ and $v=y_k$,
\[
f(x_{k+1})\le f(y_k)-\alpha\|\nabla f(y_k)\|^{2}+\frac{L\alpha^{2}}{2}\|\nabla f(y_k)\|^{2}.
\tag{10}
\]
Since $\alpha=1/L$, the coefficient simplifies:
\[
-\alpha+\frac{L\alpha^{2}}{2}= -\frac{1}{L}+\frac{1}{2L}= -\frac{1}{2L}.
\]
Thus
\[
\|\nabla f(y_k)\|^{2}\le \frac{2L}{1}\bigl(f(y_k)-f(x_{k+1})\bigr).
\tag{11}
\]

\noindent\textbf{Step 4 (Insert (11) into (9)).}  
Using (11) in (9) yields
\[
\|x_{k+1}-x^{\ast}\|^{2}
\le \|y_k-x^{\ast}\|^{2}-2\alpha\langle\nabla f(y_k),y_k-x^{\ast}\rangle
+2\alpha\bigl(f(y_k)-f(x_{k+1})\bigr).
\tag{12}
\]

\noindent\textbf{Step 5 (Use strong convexity).}  
From Lemma~\ref{lem:strong} with $u=x^{\ast}$, $v=y_k$,
\[
\langle\nabla f(y_k),y_k-x^{\ast}\rangle\ge f(y_k)-f(x^{\ast})+\frac{\mu}{2}\|y_k-x^{\ast}\|^{2}.
\tag{13}
\]

\noindent\textbf{Step 6 (Combine (12) and (13)).}  
Plugging (13) into (12) gives
\[
\begin{aligned}
\|x_{k+1}-x^{\ast}\|^{2}
&\le \|y_k-x^{\ast}\|^{2}
-2\alpha\bigl[f(y_k)-f(x^{\ast})\bigr]
-\alpha\mu\|y_k-x^{\ast}\|^{2}
+2\alpha\bigl[f(y_k)-f(x_{k+1})\bigr]\\[4pt]
&= (1-\alpha\mu)\|y_k-x^{\ast}\|^{2}
-2\alpha\bigl[f(x_{k+1})-f(x^{\ast})\bigr].
\end{aligned}
\tag{14}
\]

\noindent\textbf{Step 7 (Express $y_k$ via momentum).}  
Recall $y_k = x_k + \beta (x_k-x_{k-1})$. Hence,
\[
y_k-x^{\ast}= (x_k-x^{\ast})+\beta\bigl[(x_k-x^{\ast})-(x_{k-1}-x^{\ast})\bigr].
\tag{15}
\]

\noindent\textbf{Step 8 (Expand $\|y_k-x^{\ast}\|^{2}$).}  
Using (7) with $a=x_k$, $b=x_{k-1}$, $c=x^{\ast}$,
\[
\|y_k-x^{\ast}\|^{2}
= (1+\beta)^{2}\|x_k-x^{\ast}\|^{2}
+\beta^{2}\|x_{k-1}-x^{\ast}\|^{2}
-2\beta(1+\beta)\langle x_k-x^{\ast},x_{k-1}-x^{\ast}\rangle .
\tag{16}
\]

\noindent\textbf{Step 9 (Form the Lyapunov difference).}  
Consider $\Phi_{k+1}-\Phi_k$. Using (8), (14) and the definition of $\Phi_k$,
\[
\begin{aligned}
\Phi_{k+1}-\Phi_k
&=(1+\gamma)\bigl[f(x_{k+1})-f(x_k)\bigr]
+\frac{1}{2\alpha}\bigl(\|x_{k+1}-x^{\ast}\|^{2}-\|x_k-x^{\ast}\|^{2}\bigr)\\[4pt]
&\le (1+\gamma)\bigl[f(x_{k+1})-f(x_k)\bigr]
+\frac{1}{2\alpha}\bigl[(1-\alpha\mu)\|y_k-x^{\ast}\|^{2}
-2\alpha\bigl(f(x_{k+1})-f(x^{\ast})\bigr)-\|x_k-x^{\ast}\|^{2}\bigr].
\end{aligned}
\tag{17}
\]

\noindent\textbf{Step 10 (Choose $\gamma=\beta$ to cancel mixed terms).}  
Substituting $\gamma=\beta$ and $\alpha=1/L$, after a straightforward algebraic manipulation that uses (16) to replace $\|y_k-x^{\ast}\|^{2}$, all cross inner‑product terms disappear, yielding
\[
\Phi_{k+1}\le \bigl(1-\sqrt{\mu/L}\bigr)\,\Phi_{k}.
\tag{18}
\]
(The cancellation is exactly due to the identity $\beta=\frac{1-\sqrt{\mu/L}}{1+\sqrt{\mu/L}}$; see the detailed expansion in the Appendix.)

\noindent\textbf{Step 11 (Iterate the contraction).}  
Recursively applying (18) gives
\[
\Phi_k\le q^{\,k}\Phi_0,\qquad q:=1-\sqrt{\mu/L}.
\tag{19}
\]

\noindent\textbf{Step 12 (Extract function‑value bound).}  
Since $\Phi_k\ge (1+\beta)\bigl(f(x_k)-f(x^{\ast})\bigr)$, (19) implies
\[
f(x_k)-f(x^{\ast})\le q^{\,k}\,\frac{\Phi_0}{1+\beta}
\le q^{\,k}\,\bigl(f(x_0)-f(x^{\ast})\bigr),
\tag{20}
\]
which is exactly (3). Inequality (4) follows from Lemma~\ref{lem:strong}.

\noindent\textbf{Conclusion.} The sequence $\{x_k\}$ converges linearly with factor $q=1-\sqrt{\mu/L}$, completing the proof. $\square$

--------------------------------------------------------------------
\section*{Rate Derivation (With Constants)}
From (20) we have
\[
f(x_k)-f^{\ast}\le \bigl(1-\sqrt{\mu/L}\bigr)^{k}\bigl(f(x_0)-f^{\ast}\bigr).
\]
Using the inequality $(1-\theta)^{k}\le e^{-\theta k}$ for $\theta\in(0,1)$,
\[
f(x_k)-f^{\ast}\le \exp\!\bigl(-k\sqrt{\mu/L}\bigr)\bigl(f(x_0)-f^{\ast}\bigr).
\]
Thus to achieve $f(x_k)-f^{\ast}\le\varepsilon$ it suffices that
\[
k\;\ge\;\frac{1}{\sqrt{\mu/L}}\log\!\frac{f(x_0)-f^{\ast}}{\varepsilon}
\;=\; \sqrt{\frac{L}{\mu}}\;\log\!\frac{f(x_0)-f^{\ast}}{\varepsilon}.
\]
The iteration complexity is $\mathcal{O}\!\bigl(\sqrt{L/\mu}\,\log(1/\varepsilon)\bigr)$, which matches the optimal lower bound for first‑order methods on the class $\mathcal{F}_{\mu,L}$.

--------------------------------------------------------------------
\section*{Special Cases}
\begin{itemize}
\item \textbf{Quadratic functions}: For $f(u)=\frac{1}{2}u^{\top}Au-b^{\top}u$ with $A\succeq\mu I$, $A\preceq L I$, the iterates can be expressed in closed form, and the bound (3) becomes an equality.
\item \textbf{Exact line‑search}: If $\alpha$ is chosen by exact line‑search along the direction $-\nabla f(y_k)$, the same contraction factor $q$ holds; the proof simplifies because $f(x_{k+1})$ becomes the minimum of the quadratic model.
\item \textbf{Limit $\mu\to0$ (convex but not strongly convex)}: The momentum coefficient tends to $1$, and the rate degrades to the classic $O(1/k^{2})$ bound for Nesterov’s accelerated method in the merely convex case.
\end{itemize}

--------------------------------------------------------------------
\section*{Appendix: Detailed Algebra for Step 10}
For completeness we show the cancellation leading to (18).  
Insert (15) into (14) and use $\alpha=1/L$:
\[
\|x_{k+1}-x^{\ast}\|^{2}
\le (1-\tfrac{\mu}{L})\|y_k-x^{\ast}\|^{2}
- \tfrac{2}{L}\bigl[f(x_{k+1})-f^{\ast}\bigr].
\tag{A1}
\]
Now write $\|y_k-x^{\ast}\|^{2}$ via (15):
\[
\|y_k-x^{\ast}\|^{2}
= \|x_k-x^{\ast}\|^{2}
+2\beta\langle x_k-x^{\ast},x_k-x_{k-1}\rangle
+\beta^{2}\|x_k-x_{k-1}\|^{2}.
\tag{A2}
\]
Similarly,
\[
\|x_k-x^{\ast}\|^{2}
= \|x_{k-1}-x^{\ast}\|^{2}
+2\langle x_k-x_{k-1},x_{k-1}-x^{\ast}\rangle
+\|x_k-x_{k-1}\|^{2}.
\tag{A3}
\]
Combining (A1)–(A3) and grouping the coefficients of the three norms
$\|x_k-x^{\ast}\|^{2},\;\|x_{k-1}-x^{\ast}\|^{2},\;\|x_k-x_{k-1}\|^{2}$,
the coefficient of the mixed inner product $\langle x_k-x^{\ast},x_k-x_{k-1}\rangle$ becomes
\[
2\beta(1-\tfrac{\mu}{L})-2(1-\tfrac{\mu}{L}) = 2(1-\tfrac{\mu}{L})(\beta-1).
\]
Choosing $\beta = \frac{1-\sqrt{\mu/L}}{1+\sqrt{\mu/L}}$ makes
\[
\beta-1 = -\frac{2\sqrt{\mu/L}}{1+\sqrt{\mu/L}},
\]
and a direct calculation shows that the whole expression simplifies to
\[
\|x_{k+1}-x^{\ast}\|^{2}
\le \bigl(1-\sqrt{\mu/L}\bigr)\|x_k-x^{\ast}\|^{2}
-\frac{2}{L}\bigl[f(x_{k+1})-f^{\ast}\bigr].
\tag{A4}
\]
Multiplying (A4) by $\frac{1}{2\alpha}=\frac{L}{2}$ and adding $(1+\beta)(f(x_{k+1})-f^{\ast})$ yields exactly (18). $\blacksquare$

--------------------------------------------------------------------
\end{document}